% main.tex -- pto-core technical manuscript.
% Compile: pdflatex main (x2). The bibliography is an inline thebibliography
% block, natbib-compatible via \citep/\citet, so no bibtex pass is required.
%
% NUMERICAL INTEGRITY. Every quantitative claim in this document resolves
% through a macro defined in numbers.tex, and every macro in numbers.tex is
% transcribed from a named artifact on disk (see that file's header comments
% and the per-macro source comments). Where no artifact supports a number for
% a claim worth making, the claim is stated as an open question in
% Section~\ref{sec:limitations} rather than estimated. Do not hand-edit a
% figure into this file; regenerate numbers.tex from the underlying report
% instead.
%
% HOUSE STYLE. This document contains no em dashes. Compound thoughts are
% joined with colons and semicolons, set off in parentheses, or split into
% separate declarative sentences. Dense mechanism is followed by a plain
% restatement ("In other words", "This means that") so the engineering claim
% lands in operational terms.

\documentclass[10pt,twocolumn]{article}

\usepackage[margin=0.75in]{geometry}
\usepackage{mathptmx}
% mathptmx remaps \rmdefault and the math fonts to Times/Symbol/Chancery
% metrics; it does not touch \ttdefault, so \ttfamily still renders as
% Computer Modern Typewriter rather than being remapped to a Times-metric
% mono face. That is deliberate here: this document is dense with \texttt{}
% for binary names, flags and identifiers, and a genuine fixed-width face is
% what makes those legible against the Times body text.
\usepackage{amsmath}
\usepackage{booktabs}
\usepackage{caption}
\usepackage[round]{natbib}
\usepackage{graphicx}
\usepackage{url}
\usepackage[colorlinks=true,allcolors=blue]{hyperref}

% numbers.tex — every quantitative macro used in main.tex.
%
% RULE: every macro here is transcribed from a specific file checked into the
% repository, named in the comment above it. Nothing is estimated, rounded up,
% or interpolated between two measurements. If a claim in main.tex has no
% macro here, it is because no artifact on disk supports a number for it, and
% the text says so in prose rather than inventing one. Re-run the underlying
% harness and update this file; never hand-edit a figure into main.tex itself.

% =====================================================================
% pto-peaks vs MACS3 concordance
% Source: reports/peaks_concordance_metrics.json,
%         pto-core/modules/peaks/CONCORDANCE.md
% Run: 2026-09-10 (re-run 2026-09-11, reproduced to the 4th decimal).
% Dataset: ENCODE ENCSR217QAB (Greenleaf lab, ENCODE4), K562 ATAC-seq,
%          replicate 2, ENCFF121ZQX.bam.
% =====================================================================
\newcommand{\EncodeAccession}{ENCSR217QAB}
\newcommand{\EncodeBamAccession}{ENCFF121ZQX}
\newcommand{\EncodeReads}{5{,}778{,}260}
\newcommand{\EncodeFragments}{2{,}889{,}130}
\newcommand{\NumMacsPeaks}{31{,}607}
\newcommand{\NumPtoPeaks}{31{,}835}
\newcommand{\NumMatchedPeaks}{18{,}434}
\newcommand{\PeaksRecall}{0.583}
\newcommand{\PeaksPrecision}{0.579}
\newcommand{\SummitOverlapRate}{0.9616}
\newcommand{\SummitOverlapPct}{96.2\%}
\newcommand{\MeanIoU}{0.458}
\newcommand{\SpearmanLogP}{0.8568}
\newcommand{\SpearmanLogQ}{0.857}
\newcommand{\ParityGate}{0.98}
\newcommand{\SummitFloor}{0.95}
\newcommand{\SpearmanFloor}{0.85}
\newcommand{\PeaksCountAgreement}{0.7\%}
% Peak width distribution (CONCORDANCE.md, "the disagreement actually is")
\newcommand{\PtoMedianWidth}{96}
\newcommand{\PtoPTen}{57}
\newcommand{\PtoPNinety}{191}
\newcommand{\MacsMedianWidth}{275}
\newcommand{\MacsPTen}{169}
\newcommand{\MacsPNinety}{547}
\newcommand{\EncodeMedianWidth}{336}
\newcommand{\WidthRatio}{0.35$\times$}
\newcommand{\ReciprocalOverlapFrac}{23.5\%}
% Matched-peak agreement (>=50% reciprocal overlap, n=7,442)
\newcommand{\NumMatchedPairsIoU}{7{,}442}
\newcommand{\MatchedSummitDist}{1}
\newcommand{\MatchedSummitWithinFifty}{0.976}
\newcommand{\MatchedSummitWithinHundred}{0.997}
\newcommand{\MatchedPearsonR}{0.990}
\newcommand{\MatchedSpearmanR}{0.908}
\newcommand{\MacsOnlyMedianLogP}{6.8}
\newcommand{\OverallMedianLogP}{11.2}
\newcommand{\PtoOnlyMedianLogP}{3.5}
\newcommand{\SplitMacsPeaks}{523}
\newcommand{\SplitMacsPct}{1.7\%}
% --extend-peaks compatibility mode
\newcommand{\ExtendPeaksDemo}{90}
\newcommand{\ExtendPeaksMedianWidth}{276}
% Loose cross-check against ENCODE's own replicated calls (ENCFF187JYM)
\newcommand{\EncodeAnyOverlapPto}{0.73}
\newcommand{\EncodeAnyOverlapMacs}{0.74}
\newcommand{\EncodeReciprocalPto}{0.03}
\newcommand{\EncodeReciprocalMacs}{0.57}

% =====================================================================
% Cross-tool Pareto benchmarks (speed / memory / correctness)
% Source: reports/TECHNICAL_BRIEF.md, reports/TECHNICAL_BRIEF.json,
%         reports/figures/investor_figures.md, reports/SUMMARY.md
% Host: darwin arm64, 12 logical CPUs. Measured 2026-08-28.
% =====================================================================
\newcommand{\NumBenchCases}{7}
\newcommand{\SpeedupMin}{4.1$\times$}
\newcommand{\SpeedupMax}{56.2$\times$}
\newcommand{\SpeedupMedian}{35.1$\times$}
\newcommand{\RssRangeLow}{0.4\%}
\newcommand{\RssRangeHigh}{101.2\%}
\newcommand{\IntermediateBytesAvoided}{2{,}090.9}

% cuttag/profile vs deeptools (bamCoverage + computeMatrix, summed)
\newcommand{\CuttagWallMedium}{2.30}
\newcommand{\CuttagRssMedium}{89.3}
\newcommand{\DeeptoolsWallMedium}{87.72}
\newcommand{\DeeptoolsRssMedium}{447.1}
\newcommand{\NumCuttagSpeedupMedium}{38.18$\times$}
\newcommand{\NumCuttagRamMedium}{5.01$\times$}
\newcommand{\CuttagWallSmall}{0.08}
\newcommand{\CuttagRssSmall}{10.0}
\newcommand{\DeeptoolsWallSmall}{4.39}
\newcommand{\DeeptoolsRssSmall}{83.6}
\newcommand{\NumCuttagSpeedupSmall}{56.27$\times$}
\newcommand{\NumCuttagRamSmall}{8.36$\times$}

% fastq/qc vs fastp
\newcommand{\FastqQcWallMedium}{0.0947}
\newcommand{\FastqQcRssMedium}{4.7}
\newcommand{\FastpWallMedium}{0.7088}
\newcommand{\FastpRssMedium}{1{,}163}
\newcommand{\NumFastqSpeedup}{7.48$\times$}
\newcommand{\NumFastqRam}{247.43$\times$}

% gtk/frip vs samtools+bedtools
\newcommand{\GtkFripWallMedium}{0.78}
\newcommand{\GtkFripRssMedium}{23.5}
\newcommand{\RefFripWallMedium}{30.24}
\newcommand{\RefFripRssMedium}{1{,}681.2}
\newcommand{\NumFripSpeedupMedium}{38.80$\times$}
\newcommand{\NumFripRamMedium}{71.54$\times$}
\newcommand{\GtkFripWallSmall}{0.08}
\newcommand{\GtkFripRssSmall}{8.7}
\newcommand{\RefFripWallSmall}{2.80}
\newcommand{\RefFripRssSmall}{198.4}
\newcommand{\NumFripSpeedupSmall}{35.11$\times$}
\newcommand{\NumFripRamSmall}{22.80$\times$}

% gtk/sizes vs samtools
\newcommand{\GtkSizesWallMedium}{0.64}
\newcommand{\GtkSizesRssMedium}{6.8}
\newcommand{\RefSizesWallMedium}{2.64}
\newcommand{\RefSizesRssMedium}{6.7}
\newcommand{\NumSizesSpeedupMedium}{4.14$\times$}
\newcommand{\GtkSizesWallSmall}{0.07}
\newcommand{\RefSizesWallSmall}{0.30}
\newcommand{\NumSizesSpeedupSmall}{4.29$\times$}

% Thread scaling, cuttag/profile [medium], 80 bins (TECHNICAL_BRIEF.md)
\newcommand{\ThreadOneWall}{16.86}
\newcommand{\ThreadFourWall}{2.91}
\newcommand{\ThreadEightWall}{2.30}
\newcommand{\ThreadFourSpeedup}{5.79$\times$}
\newcommand{\ThreadEightSpeedup}{7.34$\times$}
\newcommand{\ThreadEightEfficiency}{92\%}
\newcommand{\ThreadOneCoresBusy}{4.32}

% Exactness of the answer under the speedup (TECHNICAL_BRIEF.md, "Concordance")
\newcommand{\FripFragmentsMedium}{3{,}999{,}989}
\newcommand{\FripInPeaksMedium}{1{,}629{,}113}
\newcommand{\FripValueMedium}{0.407279}
\newcommand{\FripFragmentsSmall}{399{,}971}
\newcommand{\FripInPeaksSmall}{148{,}680}
\newcommand{\FripValueSmall}{0.371727}
\newcommand{\FastqReadsExact}{1{,}168{,}663}
\newcommand{\FastqBasesExact}{29{,}216{,}575}
\newcommand{\ProfilePearsonMedium}{0.999977}
\newcommand{\ProfilePearsonSmall}{0.999975}
\newcommand{\SizesMedianLenMedium}{174}
\newcommand{\SizesMedianLenSmall}{173}
\newcommand{\SizesKsStat}{0.0}

% genomic_toolkit README.md — the sort-pipeline comparison (demo.bam, a
% different measurement context than the reports/ harness above: full
% samtools-sort+bedtools-bamtobed+sort+intersect+wc chain vs frip --keep-dups)
\newcommand{\GtkDemoRefWall}{1.42}
\newcommand{\GtkDemoPtoWall}{0.05}
\newcommand{\GtkDemoSpeedup}{28$\times$}
\newcommand{\GtkDemoFragmentsInPeaks}{148{,}680}

% =====================================================================
% fastq_stream architecture measurements
% Source: pto-core/modules/fastq_stream/README.md ("Measured performance",
% "Input framing matters more than anything else"). Apple M-series, 12 cores,
% NEON, 150 bp reads, best-of-N on a hybrid-core laptop (see caveat in text).
% =====================================================================
\newcommand{\FqBgzfThroughput}{914}
\newcommand{\FqGzipThroughput}{372}
\newcommand{\FqFramingSpeedup}{2.5$\times$}
\newcommand{\FqSpscLatency}{14}
\newcommand{\FqSpscRate}{70}
\newcommand{\FqPeakRss}{10}
\newcommand{\FqMaxThroughput}{4{,}304}
\newcommand{\FqRecordLimit}{64}
\newcommand{\FqAdapterOverlapDefault}{4}
\newcommand{\FqAdapterOverlapRec}{8}
\newcommand{\FqAdapterSpuriousDefault}{2.2\%}
\newcommand{\FqAdapterSpuriousRec}{0.15\%}

% =====================================================================
% genomic_toolkit / cuttag_profiler / scrna_matrix audit findings
% Source: pto-core/docs/AUDIT_2026-09-11_*.md, TORTURE_2026-08-17.md,
% TORTURE_2026-09-10.md, REVIEW_2026-08-15.md
% =====================================================================
% Overflow reachability thresholds (AUDIT_2026-09-11_genomic_toolkit.md, G6)
\newcommand{\GtkOverflowLinesSigned}{4{,}612}
\newcommand{\GtkOverflowLinesUnsigned}{9{,}224}
\newcommand{\GtkOverflowFragmentsSigned}{4.7\times10^{9}}
\newcommand{\GtkOverflowFragmentsUnsigned}{9.3\times10^{9}}
% cuttag_profiler C3 bin-geometry ceiling
\newcommand{\CuttagMaxWindowBp}{922}
\newcommand{\CuttagMaxBins}{10{,}000{,}000}
% scrna_matrix SM1 overflow example and SM10 capacity ceiling
\newcommand{\ScrnaOverflowK}{2^{62}}
\newcommand{\ScrnaCapacityCeiling}{2{,}147{,}483{,}647}
% Phred+64 misdetection (AUDIT_2026-09-11_fastq_stream.md, A1)
\newcommand{\PhredMisreadMeanQThirtyThree}{19.85}
\newcommand{\PhredMisreadMeanQSixtyFour}{50.85}
\newcommand{\PhredMisreadQThirtyRateLow}{0.265}
\newcommand{\PhredMisreadQThirtyRateHigh}{1.000}
% HTTP request-body DoS amplification (REVIEW_2026-08-15.md, A2)
\newcommand{\HttpDosBeforeMb}{6.2}
\newcommand{\HttpDosAfterGb}{1.51}
\newcommand{\HttpDosAmplification}{7.5$\times$}
% Float dot-product accumulator cliff (REVIEW_2026-08-15.md, A3)
\newcommand{\FloatCliffRowsWrong}{300}
\newcommand{\FloatCliffRowsTotal}{300}
\newcommand{\FloatCliffErrorBelow}{7.9\times10^{-8}}
\newcommand{\FloatCliffErrorAbove}{1.1\times10^{-1}}
% fastq_stream mutation fuzz corpus (AUDIT_2026-09-11_fastq_stream.md)
\newcommand{\FqFuzzMutants}{900}
% Findings counts, by document (confirmed-and-fixed defects only; the one
% false positive in TORTURE_2026-08-17.md, T6, is excluded and discussed in
% prose; X1 in AUDIT_2026-09-11_fastq_stream.md is a pto-cloud finding,
% outside pto-core's scope, and is likewise excluded from these counts).
\newcommand{\NumReviewFindings}{5}
\newcommand{\NumReviewFindingsRoundTwo}{3}
\newcommand{\NumReviewFindingsRoundThree}{4}
\newcommand{\NumTortureAugTwo}{5}
\newcommand{\NumTortureSepTenFastq}{14}
\newcommand{\NumTortureSepTenPeaks}{8}
\newcommand{\NumTortureSepTenGtk}{1}
\newcommand{\NumAuditFastq}{8}
\newcommand{\NumAuditGtk}{9}
\newcommand{\NumAuditCuttag}{6}
\newcommand{\NumAuditScrna}{11}
\newcommand{\NumAuditTotal}{74}
% Per-module totals, summed across every pass that touched that module
% (REVIEW_2026-08-15.md's three rounds, TORTURE_2026-08-17.md,
% TORTURE_2026-09-10.md's two passes, and that module's own
% AUDIT_2026-09-11_*.md). Cross-checked to sum to \NumAuditTotal{}.
\newcommand{\NumDefectsFastqTotal}{27}
\newcommand{\NumDefectsCuttagTotal}{9}
\newcommand{\NumDefectsScrnaTotal}{18}
\newcommand{\NumDefectsGtkTotal}{12}
\newcommand{\NumDefectsPeaksTotal}{8}
\newcommand{\NumDefectsFastqHighCrit}{9}
\newcommand{\NumDefectsCuttagHighCrit}{5}
\newcommand{\NumDefectsScrnaHighCrit}{4}
\newcommand{\NumDefectsGtkHighCrit}{4}
\newcommand{\NumDefectsPeaksHighCrit}{2}

% =====================================================================
% scrna_matrix — HNSW at 1.3M cells
% Source: pto-core/modules/scrna_matrix/docs/BENCHMARKS.md
% Hardware: Apple M4 Pro, 12 cores, Release + OpenMP, scalar SIMD kernel
% (arm64 has no AVX2/AVX-512; see Limitations).
% =====================================================================
\newcommand{\ScrnaCells}{1.3}
\newcommand{\ScrnaHnswWall}{37.1}
\newcommand{\ScrnaHnswRss}{1.01}
\newcommand{\ScrnaExtrapolatedHours}{1.6}
\newcommand{\ScrnaProjectedSpeedup}{155$\times$}
\newcommand{\ScrnaExactNsPerNSq}{3.5}
% HNSW on a 50-dim clustered embedding (the regime it is for)
\newcommand{\ScrnaHnswFiftyKSpeedup}{11.2$\times$}
\newcommand{\ScrnaHnswFiftyKRecall}{0.997}
% HNSW on raw 1,000-dim gene space (the regime it is not for)
\newcommand{\ScrnaHnswGeneSpaceSlowdown}{0.05$\times$}
\newcommand{\ScrnaHnswPcaSpeedupTenK}{2.4$\times$}

% =====================================================================
% Live deployment validation
% Source: reports/SUMMARY.md (battle test), pto-cloud/ARCHITECTURE.md
% =====================================================================
\newcommand{\LiveJobId}{61035edb-4b23-4d64-a066-3d94799c0ac8}
\newcommand{\LiveJobLatency}{27.2}


\setlength{\parskip}{0pt}

\title{\vspace{-1em}\textbf{A Hardened, Memory-Bounded C++20 Engine for\\
       High-Throughput Omics}}
\author{Christopher Alexander Perez}
\date{}

\begin{document}
\maketitle

\begin{abstract}
\noindent
Genomics preprocessing and epigenomic analysis run on tool chains that
decompress, materialize and re-read an intermediate file at nearly every
stage boundary, and on interpreted wrappers whose failure modes are inherited
from the C libraries beneath them rather than audited. I present
\texttt{pto-core}, five independent C++20 engines: a streaming Poisson peak
caller, a memory-bounded fragment-QC toolkit, a CUT\&Tag/CUT\&RUN signal
profiler, a lock-free FASTQ streaming engine, and a cache-aligned
sparse-matrix and approximate-$k$NN library for single-cell RNA-seq. All five
are built on three non-negotiable properties: bounded heap execution
independent of input size, fail-closed parsing (a corrupt record is a
refusal, never a plausible default), and correctness demonstrated under
adversarial sanitizer verification rather than assumed from a green test
suite.

Across \NumBenchCases{} workloads measured against the tools they replace
(samtools, bedtools, deeptools, fastp), \texttt{pto-core} is
\SpeedupMin{} to \SpeedupMax{} faster (median \SpeedupMedian{}) at
\RssRangeLow{} to \RssRangeHigh{} of the reference tool's peak resident
memory, and every measured answer agrees with the reference either exactly
or within a stated, physically motivated tolerance
(Table~\ref{tab:pareto-benchmarks}, Table~\ref{tab:concordance-exactness}).
Seven adversarial verification passes found and fixed \NumAuditTotal{}
confirmed defects across the five modules (Table~\ref{tab:audit-summary}).
Most of those were not crashes; they were silent-wrong-answer and
denial-of-service classes, each closed with a regression test and a
sanitizer-clean re-verification.

I report this work with the same discipline I built into it, which means
stating two results as split decisions rather than rounding them into
successes. The streaming peak caller (\texttt{pto-peaks}) reproduces MACS3's
summit locations to a \MatchedSummitDist{}\,bp median offset and its
peak-strength ranking at Pearson $r = \MatchedPearsonR{}$, but it calls
systematically narrower intervals (median \PtoMedianWidth{}\,bp against
MACS3's \MacsMedianWidth{}\,bp), so the pre-registered \ParityGate{}
interval-overlap parity gate is \textbf{not met} (mean IoU \MeanIoU{}). That
gap is documented here as a boundary-convention difference; it is not
presented as a passing regression test. The sparse-matrix engine's HNSW
backend processes \ScrnaCells{} million cells in \ScrnaHnswWall{}\,s at
\ScrnaHnswRss{}\,GB, closing a feasibility gap against an extrapolated
\ScrnaExtrapolatedHours{}-hour exact computation, but the comparative claims
against \texttt{scanpy.pp.neighbors()} that motivated the module remain
explicitly unmeasured. Source, build instructions and every report cited here
are available under the MIT license at
\url{https://github.com/chrisaperez/pto-core}.
\end{abstract}

% ======================================================================
\section{Introduction}
% ======================================================================

\subsection{Motivation}

I built \texttt{pto-core} out of a specific frustration, and I state it
plainly because it shaped every architectural decision that follows rather
than serving as a preface to them. A large fraction of production
bioinformatics software is a thin interpreted orchestration layer (Python or
R) wrapped around a C or C++ core that the orchestration layer's author did
not write and frequently has not read, invoked through a shell pipeline that
materializes an intermediate file at every stage boundary. The observable
properties that decide whether such a tool survives contact with real data,
namely its memory footprint, its behavior on malformed input and its thread
scaling, are therefore not designed at all. They are inherited from whichever
library happened to be linked. This means that the most important performance
and correctness characteristics of the software are accidental properties of
its dependency graph. That is not a tuning problem: it is an architectural
one.

Three consequences of that regime recur, and this project set out to refuse
each of them structurally rather than patch them case by case. The first is
memory bloat that scales with input size when nothing in the algorithm
requires it, typically because an intermediate file is fully materialized
where a stream would do. The second is parsers that fail \emph{open}:
a corrupted field becomes a plausible default, a truncated file becomes a
short but valid one, and a partially read stream is treated as a clean
end-of-file. The third is pipelines that fail \emph{mid-run} and
non-deterministically, in a way a retry masks rather than explains. The
second of these is the one that should alarm a working scientist most,
because its output is not an error. Its output is a number, and the number is
wrong.

I want to be careful about the scope of that critique. None of it is a claim
about any specific named tool's internals, and I make no such claim without a
reproduced artifact to back it; that is the same standard this document holds
its own numbers to (Section~\ref{sec:results}). The claim is about a failure
\emph{class}, and the best evidence I can offer for its reality is
adversarial evidence against my own work. Building five engines
\emph{deliberately} against these three patterns still produced
\NumAuditTotal{} confirmed instances of them before sustained sanitizer-driven
verification found and closed each one (Section~\ref{sec:verification}). The
lesson I draw is not that any particular library is unusually bad. It is that
this class of defect is the default outcome of writing performance-path C++
without an adversarial verification discipline applied to it continuously,
and that the discipline, not the raw throughput, is the harder and more
transferable half of this work.

\subsection{A convergence of interests}

The position I take here sits at an intersection whose components are
individually common and jointly rare: hardware-aware systems programming
(cache-line layout, runtime SIMD dispatch across NEON, AVX2 and AVX-512BW,
lock-free single-producer/single-consumer queues); a verification discipline
closer to formal-methods practice than to conventional unit testing
(sanitizer-driven fuzzing against explicit fail-closed invariants,
differential testing against arbitrary-precision references); and
computational biology's own domain constraints (BGZF block framing, 0-based
half-open coordinates, and the FRiP and Benjamini--Hochberg statistics by
which a result is judged). Each of those disciplines has a mature literature.
Very little software sits in all three at once.

I propose that the correct response to the status quo is to stop treating
bioinformatics tooling as disposable scripting and start treating it as
safety-critical systems engineering. The justification is not aesthetic. A
peak caller's output becomes a figure, the figure becomes a claim, and the
claim outlives every detail of the pipeline that produced it. Software whose
silent failure mode is a plausible wrong number carries the same class of
risk as an unaudited control system; the difference is only how long it takes
for the failure to surface. Each of the five modules below is small enough to
read end to end, and that is a deliberate design constraint rather than a
happy accident: a correctness claim about a sequencing-scale reduction is
only as credible as the code making the claim is auditable.

\subsection{Design ethos and contributions}

Three properties hold across all five modules, not as aspirations but as
invariants with regression tests behind them:

\begin{itemize}\itemsep2pt
\item \textbf{Bounded heap execution.} Peak resident memory is a property of
  the pipeline topology (buffer-pool size, ring capacity, batch size) and not
  of input size. \texttt{fastq\_stream}'s peak RSS is \FqPeakRss{}\,MB
  independent of file size (\S\ref{sec:fastq-stream});
  \texttt{genomic\_toolkit} processes fragments in 4096-record,
  16-byte-per-record batches with no sort and no intermediate file
  (\S\ref{sec:gtk}).
\item \textbf{Fail-closed parsing.} A malformed record is a refusal, not a
  default. The rule is repository-wide: a corrupt MAPQ column must not
  silently pass a quality filter, and a truncated BGZF block must not
  silently report a shorter library as a complete one. It is also the single
  most frequently rediscovered defect class in the audit trail
  (Section~\ref{sec:verification}); the same failure to check a parse result
  was found and fixed independently in \emph{four} different modules across
  three separate review passes.
\item \textbf{Verification over assertion.} Every module ships a
  dependency-free test suite (no GoogleTest, no Catch2, because these tools
  must build on a bare cluster that has a compiler and nothing else) and has
  been driven under AddressSanitizer~\citep{serebryany2012asan} and
  UndefinedBehaviorSanitizer with \texttt{-fno-sanitize-recover=all}, under
  mutation fuzzing, and, where concurrency is load-bearing, under
  ThreadSanitizer~\citep{serebryany2009tsan}.
\end{itemize}

The contributions, each backed by a citable artifact in the repository, are
as follows:

\begin{enumerate}\itemsep2pt
\item Five independently buildable C++20 engines covering streaming peak
  calling, fragment QC, chromatin signal profiling, FASTQ preprocessing and
  single-cell sparse linear algebra, unified by a deployment constraint (run
  on the host that already holds the data) rather than by a shared code
  layer.
\item A cross-tool benchmark harness that reports speed, peak memory
  \emph{and} answer agreement together, so that a speedup is never reported
  over a different answer (Section~\ref{sec:results}).
\item An adversarial verification record covering \NumAuditTotal{} confirmed
  defects across seven passes (Table~\ref{tab:audit-summary}), maintained as
  a living document rather than a one-time disclosure. It includes four cases
  in which an already-fixed defect class recurred in a sibling code path,
  which I analyze as a methodological finding in its own right
  (Section~\ref{sec:verification}).
\item A quantified concordance study against the field's incumbent peak
  caller that reports a pre-registered parity gate as \textbf{not met}, with
  the measured shape of the disagreement, rather than redefining the gate
  after the fact in order to pass it (Section~\ref{sec:peak-concordance}).
\end{enumerate}

To see how these principles interact, we begin by defining the five engines
and the deployment constraint that unifies them, then work through the
mechanism of each in turn (Section~\ref{sec:architecture}). Next, we examine
the verification methodology that produced the defect record and the
invariants that now prevent its recurrence (Section~\ref{sec:verification}).
We then evaluate the engines against the tools they replace, including the
one concordance study whose pre-registered gate was missed
(Section~\ref{sec:results}). Finally, I set out what this system does not yet
establish (Section~\ref{sec:limitations}).

% ======================================================================
\section{System Overview}
% ======================================================================

\begin{table}[t]
\centering
\caption{The five \texttt{pto-core} modules. Each is a self-contained CMake
project with its own tests and install rules; there is no shared runtime or
code layer between them by design.}
\label{tab:modules}
\small
\begin{tabular}{ll}
\toprule
Module & Role \\
\midrule
\texttt{pto-peaks} & Streaming Poisson peak calling \\
\texttt{genomic\_toolkit} & Fragment size/dup/FRiP QC \\
\texttt{cuttag\_profiler} & CUT\&Tag/CUT\&RUN signal matrices \\
\texttt{fastq\_stream} & In-memory FASTQ QC and trimming \\
\texttt{scrna\_matrix} & Sparse CSR + exact/HNSW $k$-NN \\
\bottomrule
\end{tabular}
\end{table}

All five modules are C++20, built with CMake, and independently buildable
with every other module deleted. What unifies them is a deployment
constraint rather than a shared abstraction: each tool is designed to run on
the machine that already holds the sequencing data, with no service
dependency and no network egress required for the computation itself. An
earlier iteration introduced a shared top-level \texttt{include}/\texttt{python}
layer across modules. It was reverted, and the repository's build
documentation records the absence of such a layer as a deliberate non-goal
rather than an oversight. Table~\ref{tab:modules} summarizes the five;
Section~\ref{sec:architecture} takes each in turn.

% ======================================================================
\section{Architecture and Systems Principles}
\label{sec:architecture}
% ======================================================================

\subsection{\texttt{pto-peaks}: streaming Poisson peak calling}
\label{sec:peaks}

\texttt{pto-peaks} calls narrow peaks from a coordinate-sorted fragment BED
or BAM in a single pass, with no sort and no intermediate bedGraph. Three
components do the work: a per-base pileup feeding three \emph{centred}
sliding background windows (1\,kb, 5\,kb and 10\,kb, held as fixed-capacity
rings rather than resized buffers), a SIMD-vectorized Poisson
tail-probability kernel, and an IDLE/IN\_PEAK/SUMMIT\_SEARCH/CLOSING state
machine that consumes the scored bases and applies Benjamini--Hochberg FDR
control across all candidates.

Two properties are architecturally load-bearing rather than incidental
optimizations. The first is that the background windows are centred and
truncated at contig ends. A trailing window would judge a peak against its
own rising edge and clip every call at its 3$'$ end; an untruncated centred
window at a contig boundary would divide by background that does not exist.
The second concerns where vectorization is possible at all. The state machine
cannot be vectorized, because the state at base $i{+}1$ depends on the state
at base $i$, but the Poisson evaluation feeding it can: bases are queued and
scored in a batch through one call into a SIMD kernel (scalar, NEON, AVX2 or
AVX-512BW, selected once at startup by runtime CPUID dispatch rather than by
a compile-time \texttt{-march} flag that would SIGILL on any machine other
than the one that built the binary). Bases that provably cannot clear the
significance cutoff never enter the queue, admitted by a closed-form bound on
the Poisson tail's median. In other words, the sequential part of the
algorithm stays sequential and only the arithmetic goes wide.

The kernel contains no libm call anywhere, and that is deliberate:
\texttt{log}, \texttt{exp} and \texttt{log1p} are open-coded from their
series expansions so that all four ISA variants have something to vectorize,
since no standard math library portably exposes a vector transcendental
entry point. Numerical degeneracy fails closed. A zero, negative, NaN or
infinite background $\lambda$ returns $p=1$ rather than the naive
evaluation's $p \to 0$. The direction of that choice is the entire point: the
naive form reports every centromere, assembly gap and unmeasured window as
the most significant peak in the genome. A caller that ranks assembly gaps
above real biology is not slightly miscalibrated. It is inverted.

BAM input is decoded directly from BGZF and alignment records by byte-shift
arithmetic rather than \texttt{reinterpret\_cast}, because BAM packs
\texttt{int32} fields at whatever offset the previous record ended on, which
leaves them routinely misaligned; the cast is undefined behavior and
\texttt{-fsanitize=alignment} reports it as such. The per-block CRC32 check is
retained rather than skipped for speed. A corrupt block that happens to
inflate into plausible alignment records is precisely the silent-wrong-answer
failure this project exists to refuse.

\subsection{\texttt{genomic\_toolkit}: memory-bounded fragment QC}
\label{sec:gtk}

\texttt{genomic\_toolkit} computes fragment size distributions, duplicate
rate and FRiP (Fraction of Reads in Peaks) in one streaming pass over a
coordinate-sorted BAM, BEDPE or fragment BED, replacing a shell pipeline that
would otherwise sort, spill to disk and make several full passes over the
data. Everything downstream of the reader speaks a 16-byte, coordinate-only
\texttt{Fragment} struct, handed to consumers in 4096-record batches
(64\,KiB, and therefore L2-resident) that are invalidated as soon as the
callback returns. Nothing is ever materialized as text.

Duplicate detection needs no global state on coordinate-sorted input. Two
fragments are duplicates when their 5$'$ ends agree, and on sorted input every
fragment sharing a start position arrives contiguously, so the only state
required is a set keyed on (end, orientation) that lives for exactly one
start position. This means that memory tracks the depth of the deepest
pile-up and never the size of the file: a 200\,GB BAM and a 200\,MB BAM have
the same resident footprint here. FRiP is then computed as a reduction
against peaks merged into a disjoint ascending cover at load time and queried
by binary search. Merging is not a performance trick. It is the definition of
the statistic: an unmerged peak file queried by naive interval intersection
double-counts a fragment spanning overlapping peaks, and can report a FRiP
above~1.

The module states its divergences rather than absorbing them. Duplicate
identity uses the \emph{aligned} fragment span where Picard
\texttt{MarkDuplicates} uses the unclipped 5$'$ position, so a pair whose two
copies were soft-clipped differently is one duplicate to Picard and two
distinct fragments here; there is no optical-duplicate distinction and no UMI
awareness, which makes \texttt{estimate\_library\_size()} a lower bound on
library complexity rather than a point estimate. A breaking change is
likewise stated rather than hidden: following an audit fix, column 5 of a
fragment BED is read as a 10x-style read-pair count only under an explicit
\texttt{--with-counts} flag. In a BED5 or BED6 that column is a score
(\texttt{bedtools bamtobed}~\citep{quinlan2010bedtools} writes MAPQ there),
and reading it unconditionally had silently multiplied every reported
statistic by the score column's value. Finally, all 64-bit accumulators
(fragment-base totals, size sums) are checked additions rather than raw
arithmetic. That guard was added after a crafted input, reachable from as few
as \GtkOverflowLinesSigned{} lines under \texttt{--max-length 0}, was measured
driving the signed accumulator past $2^{63}$
(Section~\ref{sec:verification}).

\subsection{\texttt{cuttag\_profiler}: reference-point signal profiling}
\label{sec:cuttag}

\texttt{cuttag\_profiler} computes CUT\&Tag and CUT\&RUN
\citep{kayaokur2019cuttag} reference-point signal matrices (loci $\times$
bins) directly from an indexed BAM in one pass. It replaces the two-stage
\texttt{deeptools} pipeline~\citep{ramirez2016deeptools}, in which
\texttt{bamCoverage} converts the BAM into a genome-wide bigWig and
\texttt{computeMatrix} then windows that track, materializing an intermediate
at genome scale even when the loci of interest cover a small fraction of the
genome. Work is distributed one locus per task across a worker pool, block
partitioned so that each worker writes a contiguous stripe of the output
matrix and the write path needs no synchronization at all. The
\texttt{htslib}~\citep{danecek2021htslib} index and header are immutable once
loaded and are shared read-only, while each worker leases its own
\texttt{htsFile} handle from a pool, because decompression state is mutable
per handle.

One deliberate anti-pattern is worth recording, since it looks like a missed
optimization. \texttt{htslib} offers a thread pool that parallelizes BGZF
block inflation, and the intuitive configuration is one pool sized to the
machine. Here that is actively harmful: with per-region parallelism already
saturating the cores, routing every block through a single shared pool
serializes inflation behind that pool's queue and adds a hand-off per block.
Rather than merely adding threads wherever a library offers them, the
profiler places parallelism at the level where the work actually divides,
which is one locus per task, and lets each handle inflate on the thread that
owns it.

Two correctness properties were closed by the 2026-09-11 audit
(\S\ref{sec:verification}) and are architecturally central rather than
incidental bug fixes. First, regions on a contig absent from the BAM are now
excluded from the column mean, from BPM's summed signal and from scaling,
while remaining in the output matrix and carrying their names. Previously
they were averaged in as a missing value of zero, which diluted the
meta-profile in direct proportion to how many regions were skipped. The
resulting curve was smooth, plausible and wrong by a constant factor the user
had no way to see. Off-contig bins inside a \emph{kept} window are still
treated as missing, which matches \texttt{deeptools}' own
\texttt{--missingDataAsZero} semantics. Second, bin geometry
(\texttt{width} $\times$ \texttt{bin\_index}) is bounded so that every product
fits in \texttt{int64\_t}: the window ceiling is \texttt{INT64\_MAX} divided
by the \CuttagMaxBins{}-bin limit, approximately
\CuttagMaxWindowBp{}\,Gbp, and it is enforced identically at the CLI and at
the \texttt{/api/profile} HTTP handler. That bound exists because an
unbounded window was found to overflow signed 64-bit arithmetic and write
wrong bin offsets in a Release build with no diagnostic whatsoever.

The interactive dashboard is an embedded HTTP server, bound to
\texttt{127.0.0.1} by default, serving assets compiled into the binary under a
strict \texttt{Content-Security-Policy}. It is structurally excluded from the
multi-tenant cloud deployment of this codebase: a
\texttt{-DPTO\_CLOUD\_BUILD=ON} build compiles the server out of the binary
entirely, verified by the absence of \texttt{socket}, \texttt{bind} and
\texttt{listen} symbols in the resulting image. The reasoning is that a
per-run session token bound to one operator's browser is the wrong isolation
primitive for a service whose isolation boundary is one ephemeral compute
task per job. In other words, the safest way to secure a listening port in
that deployment is to ensure it cannot be compiled into existence.

\subsection{\texttt{fastq\_stream}: lock-free in-memory FASTQ streaming}
\label{sec:fastq-stream}

\texttt{fastq\_stream} performs single-pass FASTQ quality control, adapter
removal and quality trimming entirely in memory, streaming clean reads
directly into an aligner through a named pipe. It eliminates the intermediate
compressed FASTQ that a conventional \texttt{zcat}, trim and \texttt{gzip}
pipeline writes and then reads back, whether the trimming stage is
Trimmomatic~\citep{bolger2014trimmomatic}, fastp~\citep{chen2018fastp} or
anything else. The engine is a five-rank pipeline (reader, an optional pool
of inflaters, a record assembler, a pool of quality-control workers, and a
writer) connected by genuinely single-producer/single-consumer lock-free ring
buffers. The word ``genuinely'' is carrying weight. An SPSC ring is correct
only with exactly one producer thread and one consumer thread, and a shared
queue drained by a worker pool is neither; so the topology is arranged such
that no ring is ever touched by more than two threads. In other words, the
queue is fast because the topology guarantees its precondition, not because a
weaker algorithm was tuned until it looked fast.

Head and tail indices are separated by
\texttt{std::hardware\_destructive\_interference\_size} so that the producer's
write and the consumer's read never contend for the same cache line, and each
side caches the opposite index in thread-private storage so a steady-state
operation never reads the other core's line at all. Measured inter-stage
hand-off latency is \FqSpscLatency{}\,ns at \FqSpscRate{}\,M items per second
(module README; see \S\ref{sec:limitations} on the provenance of these
single-host figures). Chunk order is preserved without a reorder buffer:
chunk $k$ is dispatched to lane $k \bmod D$, each lane preserves order
internally, and a round-robin collector recovers stream order for free. This
holds because a worker emits exactly one output buffer per input chunk, which
trimming guarantees since it only ever shrinks records.

The most consequential architectural finding in this module is that
\emph{input framing, not the QC arithmetic, determines whether decompression
can be parallelized at all}. \texttt{libdeflate}~\citep{biggers2016libdeflate}
exposes no streaming interface; it inflates a complete gzip member into a
caller-sized buffer, which is the source of its speed rather than an API
oversight. A single-member plain gzip stream, which is what
\texttt{bcl2fastq}, DRAGEN and \texttt{gzip} itself produce, has exactly one
member spanning the whole file, so no interior DEFLATE block boundary can be
located without inflating everything before it. Parallel decompression of
such a stream is therefore structurally impossible, for any tool, no matter
how many cores it is given. BGZF, the SAM/BAM block-gzip framing, is instead
a sequence of independent members that each declare their own compressed size
and expand to at most 64\,KiB, so members can be located by header arithmetic
and inflated concurrently. Measured on this engine, BGZF sustains
\FqBgzfThroughput{}\,MB/s of compressed throughput against
\FqGzipThroughput{}\,MB/s for plain gzip on identical content. That
\FqFramingSpeedup{} difference costs nothing to obtain: it follows from
running \texttt{bgzip} rather than \texttt{gzip} once, at the moment the
FASTQ is generated, and it then benefits every downstream consumer of that
file permanently. The bottleneck was never the arithmetic. It was the
container.

Peak resident memory is structurally bounded by
$(\text{lanes} \times \text{slots} + \text{in-flight}) \times 256\,\text{KiB}$
and is independent of input size, measured at \FqPeakRss{}\,MB on a 200k-read
run. Quality kernels (per-cycle sums, $Q{\geq}20$ and $Q{\geq}30$ fractions,
adapter seed-and-extend) are runtime-dispatched across scalar, NEON, AVX2 and
AVX-512BW paths, and each is asserted bit-identical to a scalar
double-precision reference at every read length from 0 to 300. That range is
chosen deliberately to cover the sub-vector tails, which is where such
kernels typically diverge from their scalar reference. A maximum record size
of \FqRecordLimit{}\,KiB (roughly 32\,kbp) is a stated capability limit
rather than an oversight; raising it means raising per-lane buffer memory,
which is a capacity decision and not a bug fix.

\subsection{\texttt{scrna\_matrix}: cache-aligned sparse matrices and
approximate \texorpdfstring{$k$}{k}-NN}
\label{sec:scrna}

\texttt{scrna\_matrix} is a header-only C++20 library built around
\texttt{Block\_CSR<T>}, a 64-byte (cache-line) aligned compressed
sparse-row/column representation with zero-copy \texttt{pybind11}
interoperability. A NumPy buffer that is already 64-byte aligned,
C-contiguous and (critically, after a defect closed in review) provably
immutable can be \emph{adopted} without a copy. The immutability requirement
is not fastidiousness: the module's AVX2 and AVX-512 gather kernels index a
dense scratch buffer with the adopted column indices using \emph{unchecked}
hardware gathers. This means that the bounds validation performed once at
construction is the only thing standing between the inner loop and an
arbitrary memory write, which is exactly why a buffer the caller can still
mutate afterward may not be adopted at all.

The module offers two $k$-nearest-neighbor backends with genuinely different
operating regimes, rather than one algorithm tuned to look good on a headline
benchmark: an exact $O(n^2)$ brute-force path that operates natively on the
sparse representation, and an approximate path over vendored
\texttt{hnswlib}~\citep{malkov2018hnsw} that must densify its input first,
since HNSW indexes only dense fixed-dimension vectors. That densification is
why the approximate backend is reported as \emph{losing} to exact brute force
by \ScrnaHnswGeneSpaceSlowdown{} in raw \mbox{1,000-dimensional} gene space:
each HNSW distance touches 1,000 dense dimensions where the exact path
touches roughly 60 non-zeros, so a smaller number of comparisons does not pay
for their individual cost. On a 50-dimensional PCA embedding with real
cluster structure the ordering reverses, and HNSW wins by
\ScrnaHnswFiftyKSpeedup{} at 50,000 cells at a recall of
\ScrnaHnswFiftyKRecall{}. In other words, the representation and not the
algorithm decides which backend wins, and the module documents that as a rule
of thumb rather than leaving each user to rediscover it by benchmarking both.

Reductions accumulate in a widened type (\texttt{WideAcc<T>}, which is
\texttt{double} for \texttt{float} input) and narrow exactly once through a
checked cast, a rule enforced at build time by a source-level scan rather
than left to code review. The reason is worth stating precisely, because the
intuition about floating-point error is wrong here. A bare \texttt{float}
accumulator does not gradually lose precision on large input; it saturates to
$\pm\infty$, and the downstream guard that maps a non-finite similarity to
zero then reports two \emph{identical} rows as each other's worst possible
match. The input was finite, no exception was raised, no diagnostic was
printed, and the neighbor graph is wrong. This exact defect class was
independently found and fixed at four separate sites across three review
passes (\S\ref{sec:verification}) before the source-scan test existed to
prevent a fifth.

Concurrent index mutation is serialized behind a mutex that covers the
label-count check and the insertion together, closing a data race that had
silently discarded half the rows of two concurrently-added sets while raising
no exception and leaving a perfectly valid index behind. Queries deliberately
do \emph{not} take that lock: \texttt{hnswlib}'s search is safe against a
concurrent insertion, and serializing queries would forfeit exactly the
parallel query throughput that releasing the Python GIL during a call exists
to provide.

% ======================================================================
\section{Formal Invariants and Adversarial Verification}
\label{sec:verification}
% ======================================================================

% table_audit_summary.tex
% Source: pto-core/docs/REVIEW_2026-08-15.md (three rounds), TORTURE_2026-08-17.md,
% TORTURE_2026-09-10.md (two passes), AUDIT_2026-09-11_fastq_stream.md,
% AUDIT_2026-09-11_genomic_toolkit.md, AUDIT_2026-09-11_cuttag_profiler.md,
% AUDIT_2026-09-11_scrna_matrix.md.
% Counts are confirmed-and-fixed defects only: a finding is counted here only
% if it was reproduced against a pre-fix binary or header before being
% patched, per every document's own stated convention. Excluded: the one
% false positive (T6, TORTURE_2026-08-17.md, a TSan report traced to an
% uninstrumented libomp) and X1 (AUDIT_2026-09-11_fastq_stream.md), a
% pto-cloud orchestration defect outside pto-core's module boundary.
\begin{table}[t]
\centering
\caption{Confirmed defects found and fixed across seven adversarial passes,
2026-08-15 to 2026-09-11, by module. ``High/Crit.'' is the subset rated High
or Critical severity in the source document.}
\label{tab:audit-summary}
\small
\begin{tabular}{lrr}
\toprule
Module & Confirmed defects & High/Crit. \\
\midrule
\texttt{fastq\_stream} & \NumDefectsFastqTotal{} & \NumDefectsFastqHighCrit{} \\
\texttt{genomic\_toolkit} & \NumDefectsGtkTotal{} & \NumDefectsGtkHighCrit{} \\
\texttt{cuttag\_profiler} & \NumDefectsCuttagTotal{} & \NumDefectsCuttagHighCrit{} \\
\texttt{scrna\_matrix} & \NumDefectsScrnaTotal{} & \NumDefectsScrnaHighCrit{} \\
\texttt{peaks} (\texttt{pto-peaks}) & \NumDefectsPeaksTotal{} & \NumDefectsPeaksHighCrit{} \\
\midrule
Total & \NumAuditTotal{} & 24 \\
\bottomrule
\end{tabular}

\vspace{4pt}
\raggedright\footnotesize
Every confirmed defect carries a regression test added in the same commit
that fixed it. Verification after each pass included, per module: a
sanitizer-clean ASan+UBSan run under \texttt{-fno-sanitize-recover=all}, the
pre-existing suite re-run to confirm no regression, and where the defect was
concurrency-related, a clean ThreadSanitizer run. \texttt{fastq\_stream}'s
2026-09-11 pass additionally drove a \FqFuzzMutants-mutant corpus (bit
flips, byte replacement, deletion, insertion, duplication and truncation
across raw, gzip and BGZF framing) through the sanitizer build at 1, 3 and 8
threads with no crash, hang or sanitizer report.
\end{table}


\subsection{Methodology}

Verification here means driving hostile input at the shipped binary or the
public header under a sanitizer build. It does not mean writing unit tests
for behavior already believed correct. The protocol is common to all five
modules and was repeated across seven distinct passes between 2026-08-15 and
2026-09-11 (Table~\ref{tab:audit-summary}). First, establish a
sanitizer-clean baseline on the pre-existing suite under AddressSanitizer and
UndefinedBehaviorSanitizer with \texttt{-fno-sanitize-recover=all}, so that
any later finding is attributable to the probe rather than to a pre-existing
gap in the harness. Second, reproduce every finding against the pre-fix
binary or header \emph{before} patching it, and record the reproduction so it
can be re-run. Third, fix the defect and add a regression test in the same
change. Fourth, re-verify the full sanitizer suite, adding ThreadSanitizer
wherever the defect was concurrency-related. Nothing is recorded as a finding
on suspicion alone.

Negative results are recorded with the same care. Where a suspected defect
could not be reproduced through the public API, the record says so
explicitly. One example is a latent finite-arithmetic overflow in a
cosine-similarity helper that \texttt{Block\_CSR}'s own input validation makes
unreachable in practice; it was fixed anyway, as defence in depth for a direct
caller of the public header, and documented as unreachable rather than
written up as an exploit that was never demonstrated.

Two purpose-built instruments recur across modules. The text-format readers
in \texttt{genomic\_toolkit}, \texttt{fastq\_stream} and elsewhere are driven
by seeded mutation fuzzers (bit flips, byte substitution, deletion,
insertion, duplication and truncation) rather than by a coverage-guided
engine. That choice is practical rather than principled:
\texttt{libFuzzer}'s runtime is not shipped with the Apple Clang toolchain
used for local development, and a harness that runs only where a
coverage-guided engine happens to be installed does not run in this project's
own CI. A fixed seed also makes every failure reproduce exactly. The Poisson
kernel in \texttt{peaks} is instead checked against arbitrary-precision
(80-digit \texttt{mpmath}) golden values rather than against itself, at grid
points chosen to stress the kernel's own stated accuracy model, which is a
difference of three terms each reaching $O(k\ln k)$ whose relative error is
machine epsilon times the largest term. The tolerances are derived from that
model; they are not observed values written down after a passing run.

\subsection{Representative findings}

The full record, with a reproduction, patch and regression test for each
finding, lives in the repository's \texttt{AUDIT}, \texttt{TORTURE} and
\texttt{REVIEW} documents. Five are worth stating here in the detail that
explains \emph{why} they were findings, because each one instantiates a
general failure mode rather than a one-off typo.

\textbf{Silent format misdetection.} \texttt{fastq\_stream} had no code path
that knew any quality encoding but Phred+33. The identical 2,000 reads,
encoded as Phred+64 and read as Phred+33, produced a mean quality of
\PhredMisreadMeanQSixtyFour{} against a correct value of
\PhredMisreadMeanQThirtyThree{}, and a $Q_{30}$ rate of
\PhredMisreadQThirtyRateHigh{} against a correct \PhredMisreadQThirtyRateLow{}.
Every read passed quality trimming untouched and the process exited 0. The
fix is a histogram-based heuristic evaluated at no per-base cost, keyed on
three properties that hold together only for Phred+64 misread as Phred+33:
nothing below Q31, nothing above Q72, and at least 1\% of bases at Q61 or
higher. An explicit \texttt{--phred-offset} override is available, and
undeclared Phred+64 input now exits with the repository-wide code meaning
``the input's shape made the answer wrong'' instead of continuing silently.

\textbf{Unchecked multiplication sizing a write.} Four of the five HNSW graph
producers in \texttt{scrna\_matrix} sized their result buffer as $n \times k$
with no overflow check, while the fifth (brute-force $k$-NN) had been patched
for exactly this in an earlier review. With $n=4$ rows and
$k=\ScrnaOverflowK{}$, that product wraps to exactly 0: both allocations
succeed against empty vectors, the search proceeds normally, and the result
write lands on a null pointer, confirmed as a store-to-null under UBSan. The
fix is a single checked helper called by every producer including brute
force, so that the definition exists exactly once. Notably, the same audit
reports two \emph{other} findings of the same shape, in which an already-fixed
defect class survived in a sibling path the earlier fix did not walk.

\textbf{Truncation that decodes cleanly.} A BGZF stream (a BAM, or a
\texttt{.gz}) cut exactly at a block boundary, which is what an interrupted
\texttt{aws s3 cp} or \texttt{cp} routinely produces, decodes without error.
The CRC of every remaining block is intact and only the trailing 28-byte
end-of-file marker is missing, so there is nothing for a record-level parser
to trip on. Both \texttt{genomic\_toolkit} and \texttt{fastq\_stream}
reported a plausible partial statistic for such a file at exit~0 before this
was closed. The fix checks for the EOF marker immediately after the header
and before any record is read, which is the rule \texttt{samtools quickcheck}
and \texttt{pto-peaks}' own BAM reader already enforced; it was applied to
the two modules that lacked it.

\textbf{A validation gap at the second entry point.}
\texttt{cuttag\_profiler}'s window and bin-count guards were added to the CLI
in \texttt{main.cpp} when the defect was first found. But
\texttt{ProfileOptions} has a second constructor in effect: the
\texttt{POST /api/profile} handler, which built the same struct from
untrusted JSON while applying only a subset of the checks. An unauthenticated
request with a sufficiently large \texttt{upstream}/\texttt{downstream} pair
sized a 640\,MB-per-region allocation and had the kernel SIGKILL the server
with no diagnostic at all, which is the exact scenario the CLI-side fix
believed it had already closed. The durable repair moved every guard onto
\texttt{ProfileOptions::validate()} itself, called by both entry points and
again by the allocating function as defence in depth. This means a future
third caller cannot reintroduce the gap by construction, rather than merely
being expected to remember the check.

\textbf{A cliff, not a slope, in a float accumulator.}
\texttt{scrna\_matrix}'s sparse dot-product kernel accumulates in
\texttt{float}. Below roughly $10^{17}$ in row magnitude, the float path is
exact enough to change nothing at all: a differential test against an exact
double reference found zero neighbor-set changes across every ranked
position. Above a measured threshold the dot product saturates to $+\infty$,
the guard mapping non-finite similarities to zero fires, and two
\emph{identical} rows are reported as maximally dissimilar. In one measured
sweep, \FloatCliffRowsWrong{} of \FloatCliffRowsTotal{} rows had a different
and wrong neighbor set once magnitudes crossed the threshold. The behavior is
not a gradual precision slope; it is a cliff, and everything past the edge is
wrong at once. The chosen fix rejects the input at the boundary, with
\texttt{Block\_CSR::validate()} bounding squared row norm by Cauchy--Schwarz
four orders of magnitude below \texttt{FLT\_MAX}, rather than widening the hot
$O(n^2)$ kernel. That is an explicit performance trade, taken deliberately
rather than as a drive-by change to the module's headline benchmark.

\subsection{What did not break}

Reporting only findings would misstate the balance of evidence, so what
survived hostile testing is recorded with the same care as what did not. The
SPSC ring-buffer memory ordering (release-store of the tail index before
publication, acquire-load before consumption) held clean under
ThreadSanitizer across every adversarial input tried.
\texttt{cuttag\_profiler} produced bit-identical output across thread counts
from 1 to 96. Zero-norm input vectors, which is what an empty cell looks
like, were measured rather than merely argued to produce no NaN anywhere in
\texttt{scrna\_matrix}: the all-zero row scores a defined similarity of 0
against everything, which is the correct degradation for a cell with no
counts, and it can never displace a real neighbor in anyone else's result.

One reported ThreadSanitizer race deserves its own note, because the
temptation to dispose of it in either direction was real. The race appeared
in \texttt{genomic\_toolkit}'s OpenMP FRiP reduction, and it was investigated
to ground rather than reflexively suppressed or accepted. It reproduced only
against an uninstrumented \texttt{libomp} runtime; it vanished when OpenMP was
disabled entirely and when the same workload ran single-threaded; and 55 runs
across 11 thread counts produced exactly one distinct result, which a
genuinely broken reduction would not do. No suppression was added. The
investigation and its four independent pieces of evidence were recorded
instead, so that a future genuine race in that file is not waved through by
an overbroad filename-based rule.

\subsection{Recurrence as a methodological finding}

The most useful observation to come out of this record is not any individual
defect but a pattern across them, and it was the second review pass that made
it explicit: of five findings in that pass, four were the \emph{same defect
class as an already-closed finding}, surviving in a sibling function, a
sibling code path, or a second entry point the earlier pass did not walk. I
hypothesized from that pass that the recurrence was structural rather than
coincidental, and the subsequent audits bore it out. A fix for an
unsigned-integer underflow in one quality statistic left a
character-for-character identical expression unfixed 50 lines below, in a
function that gates a filtering decision rather than a printed report. A
float-square overflow fix applied to one norm computation left three
hand-rolled sibling loops untouched. The HTTP validation gap described above
is the same pattern applied to an entry point rather than to a function.

The cause is visible once named: the audit discipline that produces good
findings, namely reproduce before fixing, also focuses the fix on the single
line where the reproduction landed. In other words, the very rigor that makes
each finding trustworthy is what narrows the repair. Three process changes
follow, and they are recorded alongside the findings themselves. After any
fix, grep the tree for the defective \emph{construct} rather than the
defective line. Validate on the type rather than at any one entry point, so
that a new caller cannot bypass a check by construction. Enumerate every
untrusted-input entry point explicitly instead of re-deriving the list from
memory during each review. I regard those three rules as at least as
transferable a contribution as any individual throughput number in
Section~\ref{sec:results}.

% ======================================================================
\section{Results}
\label{sec:results}
% ======================================================================

To evaluate these claims, we begin by defining how every figure below was
obtained. All wall-clock and memory measurements come from a single host
(darwin arm64, 12 logical CPUs) through a harness that runs the reference
tool and \texttt{pto-core} in separate processes, so that
\texttt{RUSAGE\_CHILDREN}'s monotonic high-water mark cannot leak one tool's
peak into the other's measurement. Peak RSS is read from
\texttt{wait4}/\texttt{ru\_maxrss} rather than from a sampler that could miss
a spike between polls. For a multi-stage reference pipeline, wall-clock time
is summed across stages while peak RSS is taken as the \emph{maximum} across
stages rather than the sum. That asymmetry is deliberate, and it runs against
this project's interest: summing RSS would inflate the baseline and flatter
\texttt{pto-core}.

\subsection{Cross-tool Pareto benchmarks}

% table_pareto_benchmarks.tex
% Source: reports/TECHNICAL_BRIEF.md, reports/figures/investor_figures.md
% Host: darwin arm64, 12 logical CPUs. Measured 2026-08-28T03:16-03:18Z.
% Wall clock is summed across stages for a multi-stage reference pipeline;
% peak RSS is the MAX across stages, never the sum (TECHNICAL_BRIEF.md,
% "How these numbers were measured") -- the one direction of error the
% harness is built to prevent from flattering pto-core.
\begin{table*}[t]
\centering
\caption{Wall-clock and peak-RSS comparison against the established tool each
workload replaces, over the full case $\times$ tier matrix that produced an
agreed answer (\NumBenchCases{} of \NumBenchCases{} cases; concordance is
Table~\ref{tab:concordance-exactness}). MB is $10^6$ bytes throughout.}
\label{tab:pareto-benchmarks}
\footnotesize
\setlength{\tabcolsep}{4pt}
\begin{tabular}{llrrlrrrr}
\toprule
Workload & Tier & pto-core (s) & Reference (s) & Reference tool & pto RSS (MB) & Ref.\ RSS (MB) & Speedup & RAM ratio \\
\midrule
\texttt{cuttag/profile} & medium & \CuttagWallMedium{} & \DeeptoolsWallMedium{} & deeptools & \CuttagRssMedium{} & \DeeptoolsRssMedium{} & \NumCuttagSpeedupMedium{} & \NumCuttagRamMedium{} \\
\texttt{cuttag/profile} & small  & \CuttagWallSmall{} & \DeeptoolsWallSmall{} & deeptools & \CuttagRssSmall{} & \DeeptoolsRssSmall{} & \NumCuttagSpeedupSmall{} & \NumCuttagRamSmall{} \\
\texttt{fastq/qc}       & medium & \FastqQcWallMedium{} & \FastpWallMedium{} & fastp & \FastqQcRssMedium{} & \FastpRssMedium{} & \NumFastqSpeedup{} & \NumFastqRam{} \\
\texttt{gtk/frip}       & medium & \GtkFripWallMedium{} & \RefFripWallMedium{} & samtools+bedtools & \GtkFripRssMedium{} & \RefFripRssMedium{} & \NumFripSpeedupMedium{} & \NumFripRamMedium{} \\
\texttt{gtk/frip}       & small  & \GtkFripWallSmall{} & \RefFripWallSmall{} & samtools+bedtools & \GtkFripRssSmall{} & \RefFripRssSmall{} & \NumFripSpeedupSmall{} & \NumFripRamSmall{} \\
\texttt{gtk/sizes}      & medium & \GtkSizesWallMedium{} & \RefSizesWallMedium{} & samtools & \GtkSizesRssMedium{} & \RefSizesRssMedium{} & \NumSizesSpeedupMedium{} & 0.99$\times$ \\
\texttt{gtk/sizes}      & small  & \GtkSizesWallSmall{} & \RefSizesWallSmall{} & samtools & 6.8 & 6.7 & \NumSizesSpeedupSmall{} & 0.99$\times$ \\
\bottomrule
\end{tabular}

\vspace{2pt}
\raggedright\footnotesize
\texttt{fastq/qc} [small] is omitted: the tier stages no \texttt{reads\_r1}
input and the case reports \texttt{skip}, not a measurement -- listed as an
omission in the underlying report rather than left out of this table
silently. Across the \NumBenchCases{} cases shown, peak RSS ranges from
\RssRangeLow{} to \RssRangeHigh{} of the reference tool's (the
\texttt{gtk/sizes} rows, where both tools stream and neither buffers,
are the two cases pto-core does not win on memory), and the reference
pipelines wrote \IntermediateBytesAvoided{}~MB of intermediate files across
these cases that pto-core never materialises at all.
\end{table*}


Table~\ref{tab:pareto-benchmarks} reports every case in which
\texttt{pto-core} and the established tool it replaces were confirmed to
agree on the answer, with the agreement evidence itself in
Table~\ref{tab:concordance-exactness} below. A speedup is shown only where
that agreement holds, which is the same discipline applied to the
peak-concordance study in Section~\ref{sec:peak-concordance}. The two
\texttt{gtk/sizes} rows are the cases where \texttt{pto-core} does not win on
memory (RAM ratio $\approx 1\times$), and the reason is instructive rather
than disappointing: both \texttt{genomic\_toolkit sizes} and
\texttt{samtools} stream the input and neither buffers it, so there is no
intermediate materialization for a streaming architecture to remove. The
\texttt{fastq/qc} [small] case is listed as skipped rather than quietly
omitted, following the harness's own convention that a report showing only
the cases that ran is indistinguishable from a report of a run in which
everything ran.

\subsection{Correctness under speed}

% table_concordance_exactness.tex
% Source: reports/TECHNICAL_BRIEF.md, "Concordance" section. The evidence that
% the speedups in Table~\ref{tab:pareto-benchmarks} are not bought with a
% different answer -- exact means bit-identical counts; a tolerance is shown
% only where the two tools compute genuinely different quantities by
% construction (e.g. deeptools' bigWig-derived normalisation vs pto-core's
% direct BAM-derived one).
\begin{table}[t]
\centering
\caption{Correctness under the speedups of Table~\ref{tab:pareto-benchmarks}.
All deltas are on the medium tier unless noted.}
\label{tab:concordance-exactness}
\footnotesize
\begin{tabular}{@{}lrrl@{}}
\toprule
Case & pto-core & Reference & Agrees \\
\midrule
\texttt{cuttag/profile}: Pearson $r$ & \ProfilePearsonMedium{} & 1.0 & tol.\ 0.01 \\
\texttt{fastq/qc}: reads & \FastqReadsExact{} & \FastqReadsExact{} & \textbf{exact} \\
\texttt{fastq/qc}: bases & \FastqBasesExact{} & \FastqBasesExact{} & \textbf{exact} \\
\texttt{gtk/frip}: fragments & \FripFragmentsMedium{} & \FripFragmentsMedium{} & \textbf{exact} \\
\texttt{gtk/frip}: in peaks & \FripInPeaksMedium{} & \FripInPeaksMedium{} & \textbf{exact} \\
\texttt{gtk/frip}: FRiP & \FripValueMedium{} & 0.4072794 & tol.\ $10^{-6}$ \\
\texttt{gtk/sizes}: median len.\ (bp) & \SizesMedianLenMedium{} & \SizesMedianLenMedium{} & \textbf{exact} \\
\texttt{gtk/sizes}: KS statistic & \SizesKsStat{} & \SizesKsStat{} & \textbf{exact} \\
\bottomrule
\end{tabular}
\end{table}


A speedup over a wrong answer is not a result, so
Table~\ref{tab:concordance-exactness} reports the same run's answer-agreement
alongside the timings above. Three of the four cases require exact,
bit-identical counts: \texttt{fastq/qc}'s read and base totals,
\texttt{gtk/frip}'s fragment and in-peak counts, and \texttt{gtk/sizes}'
median length together with the Kolmogorov--Smirnov statistic against the
reference size distribution. \texttt{cuttag/profile}'s meta-profile is
compared by Pearson correlation under a stated tolerance instead, because
\texttt{deeptools} and \texttt{cuttag\_profiler} normalize against different
denominators by construction (a bigWig-derived coverage track versus fragment
counts read directly from the BAM). This means the property under test there
is the \emph{shape} of the profile, not numerical identity, and the tolerance
says so rather than hiding a mismatch.

\subsection{Thread scaling and the cores-busy caveat}

On the \texttt{cuttag/profile} [medium] workload (80 bins), wall-clock time
falls from \ThreadOneWall{}\,s at one thread to \ThreadFourWall{}\,s at four
(\ThreadFourSpeedup{}) and \ThreadEightWall{}\,s at eight
(\ThreadEightSpeedup{}, \ThreadEightEfficiency{} parallel efficiency). That
figure needs a caveat the harness itself surfaces: the nominal one-thread run
already averaged \ThreadOneCoresBusy{} cores busy, because \texttt{htslib}
decompresses BGZF on its own internal threads regardless of the
\texttt{--threads} flag passed to the profiler. In other words, the baseline
was never serial, so the efficiency denominator is wrong in this project's
favor, and any figure that reads as superlinear scaling is a property of that
baseline rather than of the algorithm. I report it this way instead of
quoting the more flattering ratio a naive one-thread-to-many-thread
comparison would produce.

\subsection{Peak-calling concordance: a split decision}
\label{sec:peak-concordance}

% table_peak_concordance.tex
% Source: reports/peaks_concordance_metrics.json,
%         pto-core/modules/peaks/CONCORDANCE.md
% Dataset: ENCODE ATAC-seq \EncodeAccession{} (\EncodeBamAccession{}.bam),
% \EncodeFragments{} fragments, GRCh38, vs MACS3 3.0.4 (-f BAMPE -q 0.05).
\begin{table}[t]
\centering
\caption{pto-peaks vs.\ MACS3 on real ENCODE K562 ATAC-seq. The pre-registered
parity bar (\ParityGate{} on both gated metrics) is stated alongside the
result it was measured against, not omitted.}
\label{tab:peak-concordance}
\footnotesize
\begin{tabular}{lrrl}
\toprule
Metric & Value & Parity gate & Met? \\
\midrule
Total MACS3 peaks & \NumMacsPeaks{} & --- & --- \\
Total pto-peaks peaks & \NumPtoPeaks{} & --- & --- \\
IoU-matched pairs & \NumMatchedPeaks{} & --- & --- \\
Summit within 50\,bp & \SummitOverlapRate{} & $\geq \ParityGate{}$ & \textbf{No} \\
Spearman $r$, $-\log_{10}p$ & \SpearmanLogP{} & $\geq \ParityGate{}$ & \textbf{No} \\
Mean IoU & \MeanIoU{} & not gated & --- \\
\bottomrule
\end{tabular}

\vspace{4pt}
\begin{tabular}{lrrr}
\toprule
Caller & Median width (bp) & p10 & p90 \\
\midrule
pto-peaks (defaults) & \PtoMedianWidth{} & \PtoPTen{} & \PtoPNinety{} \\
MACS3 & \MacsMedianWidth{} & \MacsPTen{} & \MacsPNinety{} \\
ENCODE (replicated, pooled) & \EncodeMedianWidth{} & --- & --- \\
\bottomrule
\end{tabular}

\vspace{4pt}
\begin{tabular}{@{}lr@{}}
\toprule
\multicolumn{2}{@{}l}{On the \NumMatchedPairsIoU{} pairs at $\geq$50\% reciprocal overlap} \\
\midrule
Median summit distance & \MatchedSummitDist{}\,bp \\
Summit within 50\,bp / 100\,bp & \MatchedSummitWithinFifty{} / \MatchedSummitWithinHundred{} \\
Pearson $r$ on peak strength & \MatchedPearsonR{} \\
Spearman $r$ on peak strength & \MatchedSpearmanR{} \\
\bottomrule
\end{tabular}
\end{table}


\texttt{pto-peaks} was benchmarked against MACS3~\citep{zhang2008macs} on
real ENCODE~\citep{encode2012} K562 ATAC-seq data (\EncodeAccession{},
\EncodeBamAccession{}.bam, \EncodeFragments{} fragments), which is the
reference result this project's own implementation plan called for. The
outcome is a genuine split decision and is reported as one.

On summit location and relative peak strength, the two callers are
effectively interchangeable. Over the \NumMatchedPairsIoU{} peak pairs at
$\geq$50\% reciprocal overlap, median summit distance is
\MatchedSummitDist{}\,bp, \MatchedSummitWithinFifty{} of summits fall within
50\,bp, and peak-strength Pearson correlation is \MatchedPearsonR{}. On
interval width they are not. \texttt{pto-peaks} calls systematically narrower
intervals, with a median of \PtoMedianWidth{}\,bp against MACS3's
\MacsMedianWidth{}\,bp, roughly \WidthRatio{} the width, and this single
structural difference is the direct cause of every failing summary metric.
Mean IoU of \MeanIoU{} sits close to the theoretical ceiling for a
\PtoMedianWidth{}\,bp call nested inside a \MacsMedianWidth{}\,bp call, and
the harness's IoU-greedy 1:1 matcher pairs only \ReciprocalOverlapFrac{} of
peaks at $\geq$50\% reciprocal overlap. This means that much of the apparent
42\% ``unmatched'' rate measures the matcher's inability to bridge a width
gap, not the two callers finding different genomic regions. Peak splitting is
minor and one-directional: \SplitMacsPeaks{} MACS3 peaks
(\SplitMacsPct{}) are covered by two or more \texttt{pto-peaks} calls, while
zero \texttt{pto-peaks} calls are covered by two or more MACS3 peaks. That
asymmetry is consistent with \texttt{pto-peaks} occasionally fragmenting one
wide MACS3 peak and never merging two.

The pre-registered parity gate for this project was \ParityGate{} on both
summit overlap and Spearman correlation of peak strength. Both are
\textbf{not met} (Table~\ref{tab:peak-concordance}). Following the 2026-09-10
result, the harness's exit code was re-baselined onto two \emph{regression
floors}, \SummitFloor{} and \SpearmanFloor{}, chosen just under the measured
\SummitOverlapRate{} and \SpearmanLogP{}. I want to be precise about what
that re-baselining does and does not mean. It changes the question the exit
code answers, from ``does this meet parity'' to ``has this drifted from its
first real-data result,'' and nothing more. The original \ParityGate{} gates
are still computed, printed and written to the metrics artifact
unconditionally, and they are still reported as not met; an exit code of 0
from the regression harness is documented, both in the repository and here,
as not a parity claim.

A compatibility mode exists and does not resolve the question either.
\texttt{--extend-peaks} pads every reported interval symmetrically without
altering summit position or significance, and at
\texttt{--extend-peaks}~\ExtendPeaksDemo{} the median width moves from
\PtoMedianWidth{}\,bp to approximately \ExtendPeaksMedianWidth{}\,bp, landing
on MACS3's own median. But the harness measures the default, unextended
configuration, and whether the \emph{default} boundary convention should
change is recorded as an open decision. Adding a flag is not the same as
answering the question.

\subsection{\texttt{scrna\_matrix} at scale: feasibility without a baseline}

On a synthetic \ScrnaCells{} million-cell dataset (1,000 genes, 6\% density),
naive extrapolation of the exact $O(n^2)$ brute-force path's measured scaling
law, approximately \ScrnaExactNsPerNSq{}\,ns per $n^2$ operation and holding
within measurement noise across a 16$\times$ range of $n$, projects
\ScrnaExtrapolatedHours{} hours of wall-clock time. The HNSW backend on the
same dataset is a real measured run rather than an extrapolation:
\ScrnaHnswWall{}\,s wall-clock at \ScrnaHnswRss{}\,GB peak RSS, including
index construction and all 1.3M queries, for a \ScrnaProjectedSpeedup{}
reduction against the exact path's extrapolated cost.

I report that as closing a \emph{feasibility} gap, in which an intractable
exact computation becomes a tractable approximate one, and explicitly not as
a validated performance claim against any competing tool. The comparative
targets that motivated this module, a 20$\times$ wall-clock speedup and a 60\%
peak-RSS reduction against \texttt{scanpy.pp.neighbors()}
\citep{wolf2018scanpy}, have not been measured at this scale, and no number
for either appears anywhere in this manuscript
(Section~\ref{sec:limitations}). Recall of the HNSW approximation is likewise
not reported above 50,000 cells, for a reason worth stating rather than
eliding: the exact baseline needed to measure recall is precisely the
computation that is infeasible at that size.

\subsection{Live deployment validation}

Beyond the single-host measurements above, the full pipeline was exercised
end to end against a deployed AWS stack rather than only against mocked AWS
services. That path comprises API submission, an SQS FIFO queue, a Lambda
dispatcher, one ephemeral Fargate task per job, and a callback consumer that
restores job state. One real job (\texttt{genomic\_toolkit frip} on an
8.5\,MB input) completed queue-to-terminal in \LiveJobLatency{}\,s. This is
evidence that the architecture in Section~\ref{sec:architecture} runs as a
deployed system; it is not a performance claim. A single job on a small input
is not a load test, and no throughput or concurrency figure is claimed from
it.

% ======================================================================
\section{Discussion}
% ======================================================================

The throughput results in Section~\ref{sec:results} are the least
interesting part of this work, and I would rather say so directly than let
them carry the argument. Most of the speedups have a structural explanation
rather than a micro-optimization one: a streaming reduction does less work
than a pipeline that materializes an intermediate at genome scale, and
binary search over a merged cover does less work than a general interval
intersection. Anyone rebuilding these tools with the same architectural
commitments should expect to land in the same neighborhood. The numbers are a
consequence of the design; they are not the contribution.

What I think does generalize is the verification record and the shape of what
it caught. Of the \NumAuditTotal{} confirmed defects, the ones that mattered
most were almost never crashes. They were a truncated file reporting half a
library at exit~0, a quality encoding misread into a plausible mean, a
skipped region diluting a meta-profile by a constant the user could not see,
and a saturating accumulator turning two identical cells into each other's
worst match. Each of those produces output that passes every downstream
sanity check a biologist would think to apply, because it is well-formed,
finite and in range. This is what makes the fail-open failure mode
categorically worse than a segfault: a crash costs an afternoon, while a
plausible wrong number can cost a paper. Fail-closed parsing is therefore not
defensive programming in the ordinary sense. It is the recognition that in
this domain, refusing to answer is a better outcome than answering
incorrectly.

The recurrence pattern in \S\ref{sec:verification} sharpens the point. Four
separate instances of an already-closed defect class surviving in a sibling
path is not a story about carelessness; it is a story about how a
reproduce-then-fix discipline naturally narrows a repair to the line where
the reproduction landed. The corrective is to move the guard onto the type,
so that the invariant holds for callers who have not been written yet. Rather
than merely fixing the defect that was found, the goal is to make the class
of defect unrepresentable at the boundary where untrusted data enters. That
is an old idea in systems and formal-methods practice. It is not yet a
default expectation in bioinformatics tooling, and I would like it to become
one.

This also frames what \texttt{pto-core} is not. It is not a claim that
existing tools are badly built by people who should have known better; most
of them were written under constraints that made different trade-offs
reasonable, and several are cited throughout this manuscript precisely
because they are the correct reference for their task. The claim is narrower
and, I think, harder to argue with: the economics have changed. A sequencing
core now generates data faster than the analysis layer's memory ceiling
scales, HPC allocations are spent re-reading files that never needed to be
written, and the cost of a silent wrong answer rises with every downstream
analysis built on it. Treating these engines as safety-critical systems
software, with bounded memory, runtime-dispatched vector kernels, fail-closed
invariants and a continuous adversarial audit, is not gold-plating under
those conditions. It is the minimum standard the data volume now demands.

% ======================================================================
\section{Limitations and Honest Residuals}
\label{sec:limitations}
% ======================================================================

I want to be transparent about the limits of this system, because each of the
following bears directly on what a reader should and should not conclude from
the numbers above.

\begin{itemize}\itemsep3pt

\item \textbf{Peak-calling interval-width parity is an open problem, not a
  solved one.} Section~\ref{sec:peak-concordance}'s \ParityGate{} gates are
  not met, and the regression floors that replaced them as the harness's exit
  criterion are explicitly a drift tripwire set after the fact rather than a
  validation bar. Three remedies are documented as options rather than
  decided: widen the default boundary convention toward MACS3's; change what
  the harness measures, weighting summit and strength agreement over
  full-interval IoU; or accept and document the divergence as a deliberate
  tighter-boundary design. None has been chosen.

\item \textbf{\texttt{scrna\_matrix}'s comparative performance claims against
  \texttt{scanpy.pp.neighbors()} are unmeasured at the scale that matters.}
  The 1.3M-cell HNSW run (Section~\ref{sec:scrna}) demonstrates feasibility,
  not superiority against a baseline. No like-for-like run against
  \texttt{scanpy} at that scale exists, and \texttt{sc.pp.neighbors} on a PCA
  representation is itself approximate and near-linear, which makes it a peer
  of the HNSW backend rather than a baseline it should categorically be
  expected to beat.

\item \textbf{\texttt{deeptools} parity for \texttt{cuttag\_profiler} was not
  re-measured after the skipped-region fix} (\S\ref{sec:cuttag}), because
  \texttt{deeptools} was not installed on the audit machine at the time of
  that fix. The fix changes the reported mean only when a region is skipped,
  and the benchmarked case in Table~\ref{tab:pareto-benchmarks} had none, so
  ``unchanged'' for that table is inferred from the logic of the fix rather
  than independently observed.

\item \textbf{AVX-512 is compiled but has never executed on hardware that
  natively supports it}, across every module that offers it
  (\texttt{fastq\_stream}, \texttt{peaks}, \texttt{scrna\_matrix}). The AVX2
  kernels were executed on x86-64 only under Rosetta 2 emulation, which runs
  AVX2 instructions but does not report their availability in CPUID; runtime
  auto-selection of the AVX2 path was therefore forced rather than observed.
  No throughput claim is made for either ISA beyond compiling and being
  numerically correct against a scalar reference.

\item \textbf{\texttt{std::bad\_alloc} resilience is untested.} macOS does not
  enforce \texttt{RLIMIT\_AS}, so no allocation failure could be induced on
  the development host for any module. A Linux runner with a hard
  address-space limit, or an allocator shim, is the stated path to closing
  this.

\item \textbf{The Python-binding sanitizer coverage is split across two hosts,
  for a specific and checkable reason.} macOS System Integrity Protection
  strips \texttt{DYLD\_*} variables from a protected process before it starts,
  and those variables are exactly what preload-based sanitizer instrumentation
  of a dynamically loaded Python extension requires; this was verified from
  inside the running interpreter rather than assumed.
  \texttt{scrna\_matrix}'s full C++ engine was verified locally under CI's
  exact sanitizer flags, but its Python binding layer's sanitizer coverage
  comes from Linux CI and not from this host.

\item \textbf{Two parser-level residuals remain open.} Plain gzip and raw
  FASTQ carry no end-of-stream marker analogous to BGZF's, so a file
  truncated exactly between records or between gzip members is, to any
  parser, a valid shorter input rather than a detectable truncation. Only an
  external checksum or manifest would catch it. Separately, no
  structure-aware fuzzer yet exists over \texttt{scrna\_matrix}'s
  \texttt{from\_scipy\_csr} boundary, which is the entry point at which one of
  the eleven confirmed findings in that module's audit (an unchecked
  narrowing conversion) was found by manual probing rather than by fuzzing.

\item \textbf{CRAM truncation is unexercised.} \texttt{genomic\_toolkit}'s
  BGZF end-of-file check does not extend to CRAM, which carries its own
  end-of-file container, and no CRAM fixture with a local reference was
  available to test against.

\end{itemize}

% ======================================================================
\section{Availability}
% ======================================================================

Source, build instructions, the full verification and benchmark record cited
throughout this manuscript, and the license are available at
\url{https://github.com/chrisaperez/pto-core} under the MIT license. All five
modules build independently with CMake $\geq$3.20 and a C++20 compiler.
\texttt{scrna\_matrix} additionally requires OpenMP, enforced as a hard
configure-time failure rather than a silent fallback to single-threaded
execution. That hard failure exists because of an earlier incident: a soft
OpenMP dependency once left \texttt{\_OPENMP} undefined, every
\texttt{\#pragma omp} compiled to serial code, and the test suite reported
100\% passed while never executing a single parallel iteration. A silently
single-threaded engine that CI signs off on is worse than a build that stops.
Every number in this manuscript resolves through \texttt{numbers.tex}, which
is checked in beside \texttt{main.tex}, so that regenerating the underlying
reports regenerates the manuscript.

% ======================================================================
\section{Conclusion}
% ======================================================================

I have presented five independent, hardware-aware C++20 engines for
high-throughput omics, unified by a deployment constraint and a verification
discipline rather than by shared code. Measured against the tools they
replace, they are consistently faster and lower in peak memory across
\NumBenchCases{} workloads without sacrificing answer correctness. But the
result I consider most transferable came from building them adversarially
rather than from the throughput figures: the same fail-open,
unbounded-memory and unvalidated-second-entry-point defect classes this
project explicitly set out to refuse were nonetheless found \NumAuditTotal{}
times before a sustained sanitizer-driven practice closed them. Writing
correct performance-path C++ for this domain is not a matter of intent. It is
a matter of continuous, reproducible, adversarial verification.

Two results here do not resolve as cleanly as a headline speedup:
peak-calling boundary-width parity, and single-cell $k$-NN performance at
scale relative to an established baseline. I have reported both as open
questions rather than rounding them into successes. The same
numerical-integrity standard that makes the benchmark numbers in
Section~\ref{sec:results} worth trusting is the standard that requires it.

% ======================================================================
\begin{thebibliography}{9}
\small

\bibitem[Bolger \textit{et~al.}(2014)]{bolger2014trimmomatic}
Bolger,~A.M., Lohse,~M. and Usadel,~B. (2014)
Trimmomatic: a flexible trimmer for Illumina sequence data.
\textit{Bioinformatics}, \textbf{30}, 2114--2120.

\bibitem[Chen \textit{et~al.}(2018)]{chen2018fastp}
Chen,~S., Zhou,~Y., Chen,~Y. and Gu,~J. (2018)
fastp: an ultra-fast all-in-one FASTQ preprocessor.
\textit{Bioinformatics}, \textbf{34}, i884--i890.

\bibitem[Danecek \textit{et~al.}(2021)]{danecek2021htslib}
Danecek,~P. \textit{et~al.} (2021)
Twelve years of SAMtools and BCFtools.
\textit{GigaScience}, \textbf{10}, giab008.

\bibitem[Biggers(2016)]{biggers2016libdeflate}
Biggers,~E. (2016) libdeflate: heavily optimized library for DEFLATE/zlib/gzip
compression and decompression. \url{https://github.com/ebiggers/libdeflate}.

\bibitem[ENCODE Project Consortium(2012)]{encode2012}
The ENCODE Project Consortium (2012)
An integrated encyclopedia of DNA elements in the human genome.
\textit{Nature}, \textbf{489}, 57--74.

\bibitem[Kaya-Okur \textit{et~al.}(2019)]{kayaokur2019cuttag}
Kaya-Okur,~H.S. \textit{et~al.} (2019)
CUT\&Tag for efficient epigenomic profiling of small samples and single cells.
\textit{Nat.\ Commun.}, \textbf{10}, 1930.

\bibitem[Malkov and Yashunin(2018)]{malkov2018hnsw}
Malkov,~Y.A. and Yashunin,~D.A. (2018)
Efficient and robust approximate nearest neighbor search using hierarchical
navigable small world graphs.
\textit{IEEE Trans.\ Pattern Anal.\ Mach.\ Intell.}, \textbf{42}, 824--836.

\bibitem[Quinlan and Hall(2010)]{quinlan2010bedtools}
Quinlan,~A.R. and Hall,~I.M. (2010)
BEDTools: a flexible suite of utilities for comparing genomic features.
\textit{Bioinformatics}, \textbf{26}, 841--842.

\bibitem[Ram\'irez \textit{et~al.}(2016)]{ramirez2016deeptools}
Ram\'irez,~F. \textit{et~al.} (2016)
deepTools2: a next generation web server for deep-sequencing data analysis.
\textit{Nucleic Acids Res.}, \textbf{44}, W160--W165.

\bibitem[Serebryany and Iskhodzhanov(2009)]{serebryany2009tsan}
Serebryany,~K. and Iskhodzhanov,~T. (2009)
ThreadSanitizer: data race detection in practice.
In \textit{Proc.\ WBIA}, pp.\ 62--71.

\bibitem[Serebryany \textit{et~al.}(2012)]{serebryany2012asan}
Serebryany,~K., Bruening,~D., Potapenko,~A. and Vyukov,~V. (2012)
AddressSanitizer: a fast address sanity checker.
In \textit{Proc.\ USENIX ATC}, pp.\ 309--318.

\bibitem[Wolf \textit{et~al.}(2018)]{wolf2018scanpy}
Wolf,~F.A., Angerer,~P. and Theis,~F.J. (2018)
SCANPY: large-scale single-cell gene expression data analysis.
\textit{Genome Biol.}, \textbf{19}, 15.

\bibitem[Zhang \textit{et~al.}(2008)]{zhang2008macs}
Zhang,~Y. \textit{et~al.} (2008)
Model-based analysis of ChIP-Seq (MACS).
\textit{Genome Biol.}, \textbf{9}, R137.

\end{thebibliography}

\end{document}
